\chapter{Introducción}
\label{cap:intro}

% Este comentario no debe aparecer en el texto extraído
Los sistemas de memoria transactiva distribuyen el conocimiento entre los miembros de un
grupo de modo que cada uno sabe quién sabe qué \cite{fuente1}. El 50\,\% de la
Zanahoria aparece aquí solo como marcador de línea.

Según el autor, «la memoria transactiva es un sistema compartido para codificar almacenar y
recuperar información en grupos pequeños» \cite[p.~3]{fuente1}.

Las redes neuronales profundas aprenden representaciones jerárquicas de los datos a partir
de grandes cantidades de ejemplos etiquetados y ajustan millones de parámetros con descenso
de gradiente estocástico.

La variable $x_i^2 + \alpha$ se define en la ecuación siguiente:
\begin{equation}
  E = m c^2 \label{eq:energia}
\end{equation}

\begin{verbatim}
def codigo_de_prueba(): return 42
\end{verbatim}

Texto con acento escrito a la antigua: canci\'on y \~nandú.
\section{Anidado}
Este archivo vuelve a incluir el capítulo y forma un ciclo.
\chapter{Introducción}
\label{cap:intro}

% Este comentario no debe aparecer en el texto extraído
Los sistemas de memoria transactiva distribuyen el conocimiento entre los miembros de un
grupo de modo que cada uno sabe quién sabe qué \cite{fuente1}. El 50\,\% de la
Zanahoria aparece aquí solo como marcador de línea.

Según el autor, «la memoria transactiva es un sistema compartido para codificar almacenar y
recuperar información en grupos pequeños» \cite[p.~3]{fuente1}.

Las redes neuronales profundas aprenden representaciones jerárquicas de los datos a partir
de grandes cantidades de ejemplos etiquetados y ajustan millones de parámetros con descenso
de gradiente estocástico.

La variable $x_i^2 + \alpha$ se define en la ecuación siguiente:
\begin{equation}
  E = m c^2 \label{eq:energia}
\end{equation}

\begin{verbatim}
def codigo_de_prueba(): return 42
\end{verbatim}

Texto con acento escrito a la antigua: canci\'on y \~nandú.
\section{Anidado}
Este archivo vuelve a incluir el capítulo y forma un ciclo.
\input{cap1}



